% Pista alg-24s-q2 · Àlgebra
% Procedència: PAU Catalunya 2024 · Convocatòria de setembre · Sèrie 3 · Qüestió 2
\documentclass[11pt,a4paper]{article}
\usepackage{pista}
\pistainfo{Catalunya 2024}{Convocatòria de setembre}{Sèrie 3}

\begin{document}

\section*{\textcolor{blauPista}{Qüestió 2}}

\noindent Considereu el sistema d'equacions següent, on $m$ és un paràmetre real:
\[
\left.\begin{array}{r}
x - 3y + mz = -2 \\
x + my + 2z = 3 \\
x + y + 2z = m
\end{array}\right\}
\]

\subsection*{\textit{a)} \normalfont Discutiu el sistema segons el valor del paràmetre $m$. \hspace*{\fill}\textnormal{[1,25~punts]}}

\pista{Escriu la matriu $A$ i la matriu ampliada $A^{'}$, i calcula el determinant d'$A$, aplicant la regla de Sarrus. Igualant-lo a zero, obtindràs una equació de segon grau amb la variable $m$.

\medskip
Llavors, discuteix el sistema com sempre: cada valor d'$m$ que anul·la el determinant d'$A$ s'estudia a part, i la resta de valors es tracten tots en un sol cas, perquè el sistema és determinat i té solució única.

\medskip
\textbf{Consell:} quan substitueixis cadascun dels valors que has trobat d'$m$, mira primer les equacions que queden. Si dues equacions tenen exactament els mateixos coeficients, compara'n el terme independent.}

\subsection*{\textit{b)} \normalfont Trobeu la solució del sistema per a $m = 0$. \hspace*{\fill}\textnormal{[0,5~punts]}}

\pista{Mira en quin dels casos de la discussió de l'apartat \textit{a)} és $m = 0$: així sabràs si el sistema té solució, i quantes en té.

\medskip
A continuació, substitueix $m = 0$ a les tres equacions. Fixa't en la 2a equació i la 3a equació: tenen els mateixos termes amb $x$ i amb $z$. Si restes la 3a equació menys la 2a equació, obtindràs directament el valor de $y$. Finalment, troba $x$ amb la 1a equació, i $z$ amb la 2a equació.}

\subsection*{\textit{c)} \normalfont Per a $m = 2$, doneu una solució $(x, y, z)$ del sistema que, a més a més, compleixi $x = 5y$. \hspace*{\fill}\textnormal{[0,75~punts]}}

\pista{Substitueix $m = 2$ a les tres equacions. Fixa't en la 2a equació i la 3a equació: tenen els mateixos termes amb $x$ i amb $z$. Si restes la 2a equació menys la 3a equació, obtindràs directament el valor de $y$.

\medskip
A continuació, la condició $x = 5y$ et dona el valor de $x$. Finalment, troba $z$ amb la 1a equació.}

\end{document}
